\section{The Multi-Armed Bandit Problem}
% Slide 4
\begin{frame}{The Multi-Armed Bandit (MAB) Problem}
\begin{itemize}
    \item Multi-Armed Bandit is spoof name for ``Many Single-Armed
          Bandits''
    \item A Multi-Armed bandit problem is a 2-tuple $(\mathcal{A}, \mathcal{R})$
    \item $\mathcal{A}$ is a known set of $m$ actions (known as ``arms'')
    \item $\mathcal{R}^a(r) = \mathbb{P}[r \mid a]$ is an unknown probability distribution over rewards
    \item At each step $t$, the AI agent (algorithm) selects an action $a_t \in \mathcal{A}$
    \item Then the environment generates a reward $r_t \sim \mathcal{R}^{a_t}$
    \item The AI agent's goal is to maximize the Cumulative Reward:
    \[
        \sum_{t=1}^{T} r_t
    \]
    \item Can we design a strategy that does well (in Expectation) for any $T$?
    \item Note that any selection strategy risks wasting time on ``duds'' while
          exploring and also risks missing untapped ``gems'' while exploiting
\end{itemize}
\end{frame}

% Slide 5
\begin{frame}{Is the MAB problem a Markov Decision Process (MDP)?}
\begin{itemize}
    \item Note that the environment doesn't have a notion of State
    \item Upon pulling an arm, the arm just samples from its distribution
    \item However, the agent might maintain a statistic of history as it's State
    \item To enable the agent to make the arm-selection (action) decision
    \item The action is then a (Policy) function of the agent's State
    \item So, agent's arm-selection strategy is basically this Policy
    \item Note that many MAB algorithms don't take this formal MDP view
    \item Instead, they rely on heuristic methods that don't aim to optimize
    \item They simply strive for ``good'' Cumulative Rewards (in Expectation)
    \item Note that even in a simple heuristic algorithm, $a_t$ is a random
          variable simply because it is a function of past (random) rewards
\end{itemize}
\end{frame}
